\subsection{Exact selection by weighted matching}
\label{sec:matching-selection}

Submodularity alone does not provide an exact polynomial maximisation
algorithm. Here a stronger property of coherent collapse does: each
move can improve a regional minimum only at its two apices, and the
reward at either global apex is bounded by that move's raw Euler gain.
This reduces the entire interaction between disjoint regions to an
ordinary weighted matching problem.

Continue with the fixed tetrahedron-disjoint family $E$ and notation
of \eqref{eq:gamma}--\eqref{eq:peeled-objective}. Put $k=|E|$,
$r_i=2g_i$, and define
\begin{equation}\label{eq:omission-rewards}
 c_{vi}=\lambda_v(\gamma_v-a_{vi})_+,
 \qquad (\gamma_v-\infty)_+:=0.
\end{equation}
An omitted move preserves its old regional minimum; $c_{vi}$ measures
the resulting saving in the vertex-link correction, relative to
performing all moves. Savings from several omitted moves at the same
vertex combine by their maximum.

\begin{lemma}[Two endpoints and a cost bound]
\label{lem:omission-reward-bound}
For each move $i$, at most two distinct global vertices have
$c_{vi}>0$, and every reward satisfies $0\le c_{vi}\le r_i$.
If $v$ is a boundary vertex, the sharper bound $c_{vi}\le g_i$ holds.
For $R=E\setminus S$, the exact gain satisfies
\begin{equation}\label{eq:complement-objective}
 \begin{split}
 F(E\setminus R)&=F(E)+H(R),\\
 H(R)&=\sum_v\max\bigl(\{c_{vi}:i\in R\}\cup\{0\}\bigr)
             -\sum_{i\in R}r_i.
 \end{split}
\end{equation}
\end{lemma}
\begin{proof}
By definition $\gamma_v\le b_{vi}$ at every incidence. Thus
\cref{lem:regional-minima} gives
\[
 c_{vi}\le\lambda_v(b_{vi}-a_{vi})\le\lambda_v g_i
          \le 2g_i=r_i.
\]
A positive reward requires $b_{vi}>a_{vi}$, so $v$ must be a global
image of an apex. There are at most two such images, including when
formal vertices are identified. At a boundary vertex $\lambda_v=1$.

Equation~\eqref{eq:minimum-subset} implies
\[
 \lambda_v\bigl(\gamma_v-m_v(E\setminus R)\bigr)
    =\max\bigl(\{c_{vi}:i\in R\}\cup\{0\}\bigr).
\]
Subtracting the raw gains of omitted moves in
\eqref{eq:peeled-objective} proves \eqref{eq:complement-objective}.
\end{proof}

Form a weighted multigraph on the global vertices. A move $i$ with
two distinct positive-reward endpoints $u,v$ supplies an edge of weight
\begin{equation}\label{eq:matching-weight}
 w_i=c_{ui}+c_{vi}-r_i.
\end{equation}
Discard edges of weight at most zero. Discard moves with fewer than
two distinct positive-reward endpoints. In particular, two formal
apices representing the same global vertex do not supply two
collectable rewards: their image is one endpoint, and its reward
occurs once in \eqref{eq:complement-objective}. Among parallel edges
retain the largest weight; ties may use the least move index. Denote
the resulting simple weighted graph by $G_E$.

\begin{theorem}[Optimal peeled subbatch by matching]
\label{thm:optimal-peeled-matching}
Let $E$ be a fixed tetrahedron-disjoint family of certified coherent
$3$--$2$ moves in a finite compact $3$-manifold triangulation whose
vertex links are spheres or discs. Then
\begin{equation}\label{eq:optimal-subset-matching}
 \max_{S\subseteq E}F(S)
   =F(E)+\max_{\text{matchings }M\subseteq G_E}\sum_{i\in M}w_i.
\end{equation}
An optimal matching gives an optimal subbatch by omitting its moves
and retaining every other move of $E$. Among equal optimum peeled
scores, a maximum-weight matching of minimum cardinality gives a
subbatch retaining the largest possible number of moves.

Consequently an exact optimum over all $2^k$ subbatches is computable
in polynomial bit complexity. For integral weights the secondary
cardinality rule is obtained by maximising
\begin{equation}\label{eq:matching-tie-weight}
 \widetilde w_i=(k+1)w_i-1.
\end{equation}
\end{theorem}

\begin{proof}
Choose a subset $R\subseteq E$ maximising $H(R)$ and, among these,
having smallest cardinality. For a selected move $i\in R$, let
$U_i$ be the set of vertices at which $i$ is the unique positive
maximiser of the selected rewards. Removing $i$ loses no reward
outside $U_i$, and loses at most $c_{vi}$ at $v\in U_i$. It also
saves its cost $r_i$.

If $|U_i|\le1$, the total lost reward is at most $r_i$ by
\cref{lem:omission-reward-bound}. Removing $i$ would therefore not
decrease $H$, contradicting the choice of $R$ with smallest
cardinality. Each move in $R$ must consequently uniquely maximise
the rewards at two distinct vertices. A move has at most two
positive-reward endpoints, so these are precisely its endpoints.
Two moves of $R$ cannot share an endpoint: they cannot both be its
unique maximiser. Therefore $R$ is a matching and
\[
 H(R)=\sum_{i\in R}(c_{ui}+c_{vi}-r_i)=\sum_{i\in R}w_i.
\]
Every one of its weights is strictly positive, since deleting a
nonpositive edge cannot decrease $H$ and reduces cardinality.

Conversely, for any matching the endpoint rewards are collected
independently, and its value of $H$ is exactly the sum of its edge
weights. Replacing a parallel edge by one of greater weight never
hurts feasibility or cardinality. Thus the preceding reductions to
$G_E$ preserve the optimum and a smallest-cardinality optimum.
Equation~\eqref{eq:complement-objective} proves
\eqref{eq:optimal-subset-matching}.

Classical maximum-weight matching is polynomial
\cite{Edmonds1965}. For the tie rule, a difference of one in integer
total weight becomes a difference of $k+1$ after scaling; any two
matching cardinalities differ by at most $k$. Maximising the
transformed total therefore first maximises the original total,
then minimises the number of edges. This adds only $O(\log(k+1))$
bits to the edge weights.
\end{proof}

The cardinality rule needs the actual transformation
\eqref{eq:matching-tie-weight}. Deleting zero-weight edges alone is
insufficient: two positive edges of weight one can tie a single
edge of weight two. The smaller matching gives the same best Euler
score while performing one more Pachner reduction.

The theorem concerns the fixed family $E$. It does not solve the
additional packing problem of choosing an optimal disjoint family
from all overlapping candidates. Also, attaining the largest peeled
Euler value is a local policy objective; it does not authenticate
an essential disc or prove that the resulting fibre has a desired
connected component.

\subsection{Interior vertices bound the number of omissions}
\label{sec:interior-omission-bound}

Let $p_E^+$ be the number of interior global vertices incident to
at least one positive reward $c_{vi}$. It is at most the total
number $p$ of interior global vertices of the triangulation.

\begin{corollary}[At most one omission per relevant interior vertex]
\label{cor:interior-omission-bound}
There exists an optimal subbatch retaining at least
\[
 |E|-p_E^+\ge |E|-p
\]
moves. The minimum-omission rule in
\cref{thm:optimal-peeled-matching} satisfies this bound. With no
interior vertices, executing all moves in $E$ is optimal. With
one relevant interior vertex, at most one move need be omitted.
\end{corollary}
\begin{proof}
If both endpoints of an edge are boundary vertices, then
$c_{ui}+c_{vi}\le g_i+g_i=r_i$, so its weight is not positive.
Every edge of $G_E$ therefore touches an interior vertex with
positive reward. Edges of a matching have disjoint endpoint sets,
so such a matching has at most $p_E^+$ edges. Apply
\cref{thm:optimal-peeled-matching}.
\end{proof}

When all vertices are boundary vertices, there is also a sequential
statement: every legal coherent $3$--$2$ move has nonnegative peeled
Euler gain. Its raw gain is $2g_i$, and the total increase of the
vertex-link correction is at most $g_i+g_i$. This remains valid
when the two apices represent one global vertex. Thus any sequence
of such moves is nondecreasing in peeled Euler characteristic.

For $M$ distinct eligible degree-three central edges, the
degree-nine conflict bound supplies a greedy tetrahedron-disjoint
family $E$ with $|E|\ge\lceil M/10\rceil$. Optimal selection within
that family retains at least
\begin{equation}\label{eq:packed-retention}
 \max\{0,\lceil M/10\rceil-p\}
\end{equation}
moves. Its total peeled gain is nonnegative, because the empty
subbatch is among the choices. A bound on the number of retained
moves and a bound on the score are distinct conclusions; neither
requires summing singleton peeled gains.

\begin{corollary}[The maintained two-pole pulling source]
\label{cor:two-pole-source}
Assume the validated one-component PD is converted to the canonical
compact exterior by the maintained \texttt{pulling} subdivision,
with its source-construction contract accepted. That triangulation
has exactly two interior global vertices. This count remains two
under the certified vertex-preserving $2$--$3$ and $3$--$2$
replacements considered here. For every fixed disjoint coherent
$3$--$2$ family $E$ along such a trajectory, an optimal subbatch
retains at least $|E|-2$ moves. Starting with $M$ candidates and
the greedy packing therefore retains at least
\[
 \max\{0,\lceil M/10\rceil-2\}.
\]
In particular, $M\ge21$ guarantees at least one retained move
at an optimum of the packed family's peeled objective.
\end{corollary}
\begin{proof}
The source construction retains the north and south poles of
the crossing-cell decomposition, each with a sphere link. Its
pulling subdivision introduces only the six directed-edge cut
points per cell, all on the truncation boundary. It introduces
no face centres or cell centres. The gluing maps preserve the
two pole labels and the cut-point locus; the source theorem
identifies exactly two global poles and a single compact torus
boundary \cite{ProveItSource}. Consequently these poles are the
only interior vertices. Certified relative-boundary moves preserve
the global vertex classes and their link types. Apply
\cref{cor:interior-omission-bound} and
\eqref{eq:packed-retention}.
\end{proof}

This corollary is specific to the maintained pulling source and
its preserved vertex classes. The legacy centred subdivision
inserts additional interior vertices and does not have the same
constant bound. Arbitrary simplification routines would also
need their own proof of the vertex-count property.

\subsection{Linear selection for the two-pole source}
\label{sec:two-pole-linear-selection}

The constant omission bound has an algorithmic consequence stronger
than applying a general blossom algorithm. Every positive edge of
$G_E$ meets an interior vertex. Thus the two poles form a vertex
cover of the positive reward graph: every edge touches this set.

\begin{theorem}[Matching with an at-most-two-vertex cover]
\label{thm:two-pole-linear-matching}
Suppose an explicit positive-weight graph is supplied together with
a vertex cover $C$ of size at most two. A maximum-weight matching
of minimum cardinality can be found using a linear number of edge
inspections, integer comparisons and arithmetic operations. For
the two-pole pulling source, this gives an $O(k)$ matching-selection
step once its reward graph has been materialized, with no external
matching backend.
\end{theorem}
\begin{proof}
First check that $C$ has at most two distinct vertices and meets
every edge. Apply the transformed weights
\eqref{eq:matching-tie-weight}. A matching has at most two edges,
because each of its edges meets $C$ and matching edges cannot
share a cover vertex.

If $|C|\le1$, the best matching is empty or consists of a single
largest-weight edge. Suppose $C=\{u,v\}$. A two-edge matching must
consist of $ux$ and $vy$, where $x,y\notin C$ and $x\ne y$.
An edge $uv$ can only occur in a singleton matching. During one
scan retain the best singleton edge, the two best edges from
$u$ to distinct vertices outside $C$, and the analogous two
edges from $v$. Parallel edges have already been compressed;
equivalently one may maintain the best edge per retained outside
endpoint during the scan. Keeping two entries at each pole
requires constant work per edge.

For any admissible pair $ux,vy$, at most one of the two retained
outside endpoints at $u$ equals $y$. Hence either $ux$ is already
retained or it can be replaced by a retained edge of no smaller
transformed weight whose outside endpoint differs from $y$.
After this replacement, apply the same argument at $v$, holding
the chosen $u$ edge fixed. Thus an optimum two-edge matching is
among the at most four combinations of the retained lists.
Compare these valid combinations against the best singleton
and the empty matching. This is constant further work after
the scan. Maximising the transformed sum gives precisely the
required primary and secondary objectives.
\end{proof}

The supplied cover is an algorithmic precondition that the
implementation checks against every positive edge. It is not
accepted as an unchecked claim about the triangulation. The
geometric wrapper obtains its candidate cover from the actual
interior vertex classes and uses this path only when at most
two relevant classes occur. Repeated endpoints, loops, parallel
edges, zero gaps and equal-weight ties are handled by the
preceding reward reduction and the same exact tie transform.

The linear bound counts operations on the already formed graph.
It is not a linear bit bound independent of coordinate size,
nor a bound for geometric preparation or collective minimum
extraction. Each transformed weight has the original coefficient
bit length plus $O(\log(k+1))$ bits. The implementation's graph
assembly uses ordinary dictionaries; the theorem's graph-scan
bound does not assert a worst-case constant bound for arbitrary
Python hash-table operations.

\subsection{Algorithmic trust and independently checkable optimality}
\label{sec:matching-certificates}

The implementation forms the reward graph using exact integers.
With a validated cover of size at most two it uses
\cref{thm:two-pole-linear-matching}, entirely in the standard
library. If the graph has no positive edges, the empty omission
set is optimal without an external matching backend. For other
inputs the general solver uses the optional NetworkX
maximum-weight-matching routine.
Its documented bound is $O(q^3)$ arithmetic operations for $q$
graph vertices, with integer-only computations when the weights
are integers \cite{NetworkX}. This is separate from the nearly
linear corner-tree maintenance bound. Integer bit lengths, graph
construction, and geometry validation must also be counted.

The returned omission set has its feasibility and attained
objective recomputed directly from
\eqref{eq:complement-objective}. This establishes its value;
maximum-weight optimality of the general result still relies
on the matching backend. The two-cover path instead uses the
explicit constant-size candidate comparison just proved. The
independent geometric replay
certifies the selected moves and their exact cocycle transport,
regardless of whether the policy found the best subset.

An optional, independent optimality checker accepts a classical
Edmonds dual witness \cite{Edmonds1965}. For each vertex take a
nonnegative rational price $y_v$, and for finitely many odd
vertex sets $B$ with $|B|\ge3$ take nonnegative rational prices
$z_B$. Check every edge $uv$ of $G_E$ against
\begin{equation}\label{eq:matching-dual-feasible}
 y_u+y_v+\sum_{B\supseteq\{u,v\}}z_B\ge w_{uv}.
\end{equation}
Then check that the proposed matching weight equals
\begin{equation}\label{eq:matching-dual-value}
 \sum_v y_v+\sum_B\frac{|B|-1}{2}\,z_B.
\end{equation}
These checks prove optimality by weak duality: a matching uses
each vertex at most once and contains at most $(|B|-1)/2$ edges
inside an odd set $B$. Summing
\eqref{eq:matching-dual-feasible} over any matching bounds its
weight by \eqref{eq:matching-dual-value}. Equality for the
proposed matching proves it is maximum. The verifier needs no
assumption that the listed odd sets are laminar and runs in
time polynomial in the explicit certificate length.

The serialized certificate uses one positive integer denominator
and nonnegative integer numerators, so these checks are exact.
The current NetworkX path does not export such a witness; the
checker is available for independently supplied certificates.
It certifies the primary score optimum, not the secondary
minimum-omission tie rule. Small instances are additionally
cross-checked by an explicitly bounded exhaustive reference
solver; this testing method is not the polynomial algorithm.

The resulting theorem removes an exponential enumeration of
subbatches from this particular coherent policy problem. It
does not bound the number of hierarchy stages, guarantee
discovery of useful coherent regions, or establish the complete
recognizer's sought quasipolynomial running time.
